\documentclass[12pt]{article}
\usepackage[utf8]{inputenc}
\usepackage[T1]{fontenc}
\usepackage[margin=1in]{geometry}
\usepackage{setspace}
\usepackage{xcolor}
\usepackage[colorlinks=true, urlcolor=blue!60!black, linkcolor=blue!60!black]{hyperref}
\setlength{\parindent}{2em}
\setlength{\parskip}{1.25em}
\renewcommand{\baselinestretch}{1.0}
\begin{document}
\begin{center}
{\LARGE \textit{West Coast Hotel Co. v. Parrish}, 300 U.S. 379 (1937)\par}
\vspace{0.8em}
{\large Case Brief --- Four-Page Edition\par}
\end{center}
\vspace{1em}

\section{Memory Jogger}

In a 5--4 decision that ended the \textit{Lochner} era of economic substantive due process, the Court overruled \textit{Adkins v.\ Children's Hospital} and upheld Washington's minimum-wage law for women as a reasonable exercise of the police power under the Fourteenth Amendment's Due Process Clause (\href{https://supreme.justia.com/cases/federal/us/300/379/}{300 U.S.\ 379}).

\section{Detailed Case Facts}

Washington's 1913 ``Minimum Wages for Women'' statute declared that the State's welfare required protecting women and minors from labor conditions harmful to their health and morals. It made it unlawful to employ women at wages inadequate for their maintenance and authorized a state committee to investigate conditions, convene conferences of employers, employees, and the public, and issue binding wage orders. Underpaid employees could sue for the full minimum notwithstanding any contract for less. The committee fixed \$14.50 for a 48-hour week for hotel chambermaids.

Elsie Parrish worked as a chambermaid at the Cascadian Hotel in Wenatchee, owned by the West Coast Hotel Company, from 1933 to 1935, always below the minimum. After her discharge she and her husband sued for \$216.19, the difference between her pay and the statutory wage. The hotel did not dispute the arithmetic but argued that the statute violated freedom of contract under \textit{Adkins} (1923).

\section{Procedural History}

The Chelan County Superior Court held the statute unconstitutional under \textit{Adkins}. The Washington Supreme Court unanimously reversed. Months earlier, in \textit{Morehead v.\ New York ex rel.\ Tipaldo} (1936), a 5--4 Court with Justice Roberts in the majority had struck down New York's similar law, noting that no one had asked it to overrule \textit{Adkins}. Here that question was squarely presented. After argument in December 1936, the Court affirmed on March 29, 1937.

\section{Judicial Votes}

Affirmed, \textbf{5--4}.
\begin{itemize}
\item \textbf{Majority:} Chief Justice Hughes, joined by Justices Brandeis, Stone, Cardozo, and Roberts.
\item \textbf{Dissenting:} Justice Sutherland, joined by Justices Van Devanter, McReynolds, and Butler (the ``Four Horsemen'').
\end{itemize}
Roberts cast his conference vote on December 19, 1936, nearly two months before President Roosevelt announced his court-packing plan, yet the decision became known as ``the switch in time that saved nine.''

\section{Holding}

Washington's minimum-wage statute for women does not violate due process, and \textit{Adkins} is overruled. The Constitution guarantees no absolute ``freedom of contract''; liberty is subject to reasonable regulation in the interests of the community, and a regulation reasonably related to a proper purpose and not arbitrary or capricious satisfies due process. A legislature may address evils piecemeal, so covering only women is not arbitrary discrimination. Hughes's broader remarks about taxpayer ``subsidy'' and deference to debatable economic judgments laid the foundation for the near-total deference later cemented in \textit{Carolene Products} and \textit{Williamson v.\ Lee Optical}.

\section{Analysis of Opinions}

\textbf{Majority (Hughes).} Because Washington's statute closely paralleled the law struck down in \textit{Adkins}, reconsideration was unavoidable. Hughes rejected \textit{Lochner}'s premise: constitutional liberty is ``liberty in a social organization,'' not an ``absolute and uncontrollable liberty.'' Citing \textit{Muller v.\ Oregon}, \textit{O'Gorman}, and \textit{Nebbia}, he stressed broad legislative discretion over employment and the rule that ``every possible presumption is in favor of'' validity. He offered two justifications: the State's interest in protecting women's health against unequal bargaining power, and, more innovatively, the recognition that sub-living wages shift costs to taxpayers, as Depression relief rolls showed, so the community need not subsidize ``unconscionable employers.'' A legislature ``is free to recognize degrees of harm'' and need not regulate every class at once.

\textbf{Dissent (Sutherland).} Rebuking reliance on Depression conditions, Sutherland insisted that constitutional meaning is fixed and that judicial restraint cannot become abdication, which would turn constitutional mandates into ``mere moral reflections.'' The statute was ``essentially arbitrary'' because it tied pay to the worker's needs rather than the value of her work, ignoring the ``just equivalence'' between pay and services. And after the Nineteenth Amendment, women should not be placed in an inferior contractual class or barred from competing for lower-paid work.

\section{Examples: Future Applications}

\begin{enumerate}
\item \textbf{Same-Side (Upheld): Gender-Neutral \$15 Minimum Wage.} A wage floor reasonably related to worker welfare and reducing public assistance easily survives deferential review.
\item \textbf{Same-Side (Upheld): Dairy Price Floor.} Building on \textit{Nebbia}, courts presume a minimum retail price for milk valid unless arbitrary.
\item \textbf{Opposite-Side (Invalid): Hiring Ban for Critics of the Governor.} \textit{West Coast Hotel} relaxed review of economic regulation, not of enumerated rights such as political speech.
\item \textbf{Opposite-Side (Invalid): Modern Women-Only Wage Law.} Under \textit{Craig v.\ Boren} and \textit{United States v.\ Virginia}, sex classifications face intermediate scrutiny, vindicating Sutherland's equality critique.
\item \textbf{Fence-Sitter: Hair-Braider Licensing.} Requiring 1,500 hours of unrelated cosmetology training usually survives, but courts of appeals have split over whether naked protectionism is ``arbitrary or capricious.''
\end{enumerate}

\section{Critique}

\textbf{Originalist and libertarian views.} Randy Barnett, Richard Epstein, and David Bernstein fault the decision for inaugurating near-total judicial abdication over economic liberty and a double standard that treats contract and property as second-class rights, leaving legislatures open to rent-seeking. They note the irony that Hughes relied on \textit{Muller}-era paternalism while Sutherland defended women's formal equality.

\textbf{Progressive and historical views.} Progressives celebrate the end of \textit{Lochner}'s anti-democratic activism and credit Hughes with exposing the false neutrality of market baselines by showing that low wages shift costs to taxpayers. Barry Cushman and other revisionists argue that Roberts's vote reflected doctrinal development, especially \textit{Nebbia}, rather than capitulation to court-packing.

\section{Key Quotations}

\begin{quote}
``What is this freedom? The Constitution does not speak of freedom of contract. It speaks of liberty and prohibits the deprivation of liberty without due process of law. In prohibiting that deprivation, the Constitution does not recognize an absolute and uncontrollable liberty. Liberty in each of its phases has its history and connotation. But the liberty safeguarded is liberty in a social organization which requires the protection of law against the evils which menace the health, safety, morals, and welfare of the people. Liberty under the Constitution is thus necessarily subject to the restraints of due process, and regulation which is reasonable in relation to its subject and is adopted in the interests of the community is due process.''
\end{quote}
\noindent --- Chief Justice Hughes, Majority Opinion, 300 U.S.\ at 391.

\begin{quote}
``The judicial function is that of interpretation; it does not include the power of amendment under the guise of interpretation. To miss the point of difference between the two is to miss all that the phrase `supreme law of the land' stands for and to convert what was intended as inescapable and enduring mandates into mere moral reflections.''
\end{quote}
\noindent --- Justice Sutherland, Dissenting Opinion, 300 U.S.\ at 404.

\begin{quote}
``The exploitation of a class of workers who are in an unequal position with respect to bargaining power and are thus relatively defenseless against the denial of a living wage is not only detrimental to their health and well being, but casts a direct burden for their support upon the community. What these workers lose in wages the taxpayers are called upon to pay. The bare cost of living must be met. \ldots\ The community is not bound to provide what is in effect a subsidy for unconscionable employers.''
\end{quote}
\noindent --- Chief Justice Hughes, Majority Opinion, 300 U.S.\ at 399.

\end{document}
